%% ========================================================================
%% chapter1.tex -- Introduction (root file, standalone compilable)
%% Topic: Promptable Medical Image Segmentation Foundation Models
%%        for Endoscopic Polyp Segmentation
%% Author: Xu, Junhan; Wu, Keyou
%% Supervisor: Dr. Donglong Chen
%% Institution: BNBU United International College
%% Format reference: [GroupID] FYP I Literature Review.docx
%% Last updated: 2026-06-11
%%
%% Notes:
%%  - This file is intended to be \input{} from lecture review.tex, but it
%%    can also be compiled standalone for chapter-level review. It carries
%%    its own preamble, identical to that of lecture review.tex.
%%  - The bibliographic references point to the shared ref.bib used by the
%%    rest of the thesis; the citation keys were chosen to match the keys
%%    already present in that file.
%% ========================================================================

\chapter{Introduction}
\label{chap:intro}
\setcounter{section}{0}

\section{Clinical motivation}
\label{sec:motivation}

Colorectal cancer (CRC) is one of the leading causes of cancer-related mortality worldwide, with GLOBOCAN 2022 estimating roughly 1.9~million new cases and 904,000 deaths~\cite{bray2024globocan}, and the early detection and removal of precancerous polyps during routine colonoscopy remains the single most effective intervention for reducing CRC mortality, with long-term follow-up of the National Polyp Study reporting a substantial reduction in CRC mortality after colonoscopic polypectomy~\cite{zauber2012nejm}. Polyp segmentation, the pixel-level delineation of polyp boundaries in an endoscopic frame, is a prerequisite for several downstream tasks: estimating the malignancy risk of a detected lesion, measuring its size for surveillance planning, and providing quantitative feedback to the operator~\cite{tudela2024frontiers}. The task is intrinsically difficult because the appearance of a polyp varies in scale, shape, and texture, the boundary between a polyp and the surrounding mucosa is often weak, and the imaging conditions differ across endoscopes, operators, and clinical sites; the resulting cross-domain shift~\cite{guan2022tbme} means that a model trained on one dataset rarely generalises to another without adaptation. Closing this gap is attractive on three fronts: at the clinical level, a promptable polyp segmentor can reduce per-frame reading time and inter-observer variability during screening colonoscopy; at the engineering level, a single reproducible evaluation pipeline lowers the cost of benchmarking future polyp-specific SAM adaptations; and at the methodological level, the polyp domain offers a controlled stress test for parameter-efficient adaptation of promptable foundation models under a small-data budget.

\section{The opportunity of promptable foundation models}
\label{sec:sam}

In recent years, the Segment Anything Model (SAM)~\cite{kirillov2023sam} has demonstrated strong zero-shot segmentation capability on natural images and has been adapted to a growing number of medical imaging domains, including CT, MRI, and ultrasound. The natural question for this thesis is therefore not whether SAM \emph{can} be applied to endoscopic polyp segmentation, but how different adaptation strategies compare under one and the same evaluation pipeline.

We address that question through a staged benchmark rather than a single-method claim. The current empirical scope includes three complementary groups: zero-shot foundation-model baselines (SAM ViT-B/L/H and SAM2 Hiera-S/Base+/Large), generic medical pre-training (MedSAM), and polyp-specific adaptation recipes (PP-SAM full-encoder fine-tuning and SAMed LoRA). This design allows the thesis to separate three effects that would otherwise be conflated: backbone capacity, generic medical pre-training, and polyp-specific adaptation.

A broader family of polyp-specific SAM adaptations proposed in 2024--2025 is surveyed in Chapter~2. The present stage treats PP-SAM as a representative full-encoder few-shot adaptation reference because it directly matches the box-perturbation setting of the project, and treats SAMed as the parameter-efficient LoRA reference needed to address the Adapter/LoRA requirement in the project brief. OP-SAM and ASPS are kept as special cases: OP-SAM is reported only as a CPG-only exploratory training-free pilot, while ASPS is reported as an external fully supervised upper-bound reference rather than as a directly ranked v3 baseline.


\section{Scope of this thesis}
\label{sec:scope}

This thesis reports the current-stage reproduction and benchmarking of SAM-family models for endoscopic polyp segmentation. It does not claim that all long-term objectives in the project brief have already been completed. Instead, the empirical work is organised as a staged study: the present stage establishes a reproducible v3 evaluation protocol, reproduces representative zero-shot and few-shot SAM-family baselines, and identifies the next alignment experiments required for a fully unified comparison.

The completed empirical core focuses on two public polyp benchmarks, Kvasir-SEG~\cite{jha2020kvasir} and CVC-ClinicDB~\cite{bernal2015cvc}. The first axis studies zero-shot foundation-model transfer by comparing SAM~\cite{kirillov2023sam} with ViT-B, ViT-L, and ViT-H image encoders, and SAM2~\cite{ravi2024sam2} with Hiera-Small, Hiera-Base+, and Hiera-Large backbones. This axis is used to examine whether larger backbones consistently improve zero-shot polyp segmentation, and whether the scaling behaviour of SAM2 follows that of SAM v1. The second axis studies polyp-specific adaptation under a fixed SAM ViT-B backbone, including vanilla SAM, MedSAM~\cite{ma2024medsam}, PP-SAM~\cite{rahman2024ppsam} at 1-shot, 5-shot, and 10-shot settings, and SAMed~\cite{zhang2023samed} LoRA as the parameter-efficient adaptation extension. At the time of this partial thesis, the PP-SAM 1/5/10-shot experiments and the SAMed LoRA 1/5/10-shot plus full-training runs have been completed and their logs and evaluation files have been archived under the formal results directory.

For all v3-aligned baselines, we use the same evaluation split and the same prompt protocol: Kvasir-SEG is evaluated on its held-out \texttt{val.txt} split of 120 images, CVC-ClinicDB is used as a cross-domain test set of 612 images, and the box prompt is generated from the ground-truth mask with a fixed $\pm30$ pixel perturbation. The primary overlap metric is the Dice coefficient~\eqref{eq:dice}; the secondary overlap metric is IoU~\eqref{eq:iou}; precision and recall quantify the over- and under-segmentation tendency; and HD95~\eqref{eq:hd95} measures boundary quality. This unified protocol is described in detail in Chapter~\ref{chap:outline}.

Two entries are reported separately from the main v3-aligned ranking. ASPS~\cite{li2024asps} is used as an external fully-supervised upper-bound reference because the available checkpoint was trained on the PraNet training split; it is therefore evaluated on the PraNet held-out test set to avoid test-on-train leakage. A fully v3-aligned ASPS retraining on Kvasir \texttt{train.txt} and evaluation on Kvasir \texttt{val.txt} plus CVC-ClinicDB is left as a next-stage experiment. OP-SAM~\cite{li2025opsam} is included only as an exploratory support-conditioned, training-free baseline. The current implementation reports the CPG-only variant; a full reproduction with CPG, SPF, and EPE is also reserved for the next stage.

Several directions in the project brief are intentionally left outside the empirical scope of this partial thesis and are listed as future work. These include a systematic prompt-quality study with multiple perturbation levels, object-scale analysis with area-stratified evaluation and scale-aware modules, explicit weak-boundary refinement beyond the HD95 metric and the structure loss used by PP-SAM, training a polyp-specific Medical SAM Adapter rather than evaluating released non-polyp adapters, extension to additional polyp benchmarks such as PolypGen~\cite{ali2024polypgen}, PICCOLO~\cite{gonzalez2022piccolo}, and SUN-SEG~\cite{sun2022sunseg}, and the 3D/video directions surveyed in Chapter~2.


\section{Research questions and hypotheses}
\label{sec:rq}

The scope above leads to five research questions (RQs). The first four are empirical questions addressed by the current-stage experiments, while the fifth is a planning question that records the remaining alignment work required by the long-term project. Each RQ is paired with a falsifiable hypothesis ($H$), and the corresponding baselines are specified to avoid a promise--execution mismatch.

\begin{itemize}
    \item \textbf{RQ1 (Foundation-model scaling):} How do SAM v1 and SAM2 backbones scale under zero-shot box-prompt polyp segmentation?
    \begin{itemize}
        \item[\textbf{$H_1$:}] Increasing backbone capacity does not monotonically improve zero-shot polyp segmentation. In particular, SAM ViT-L is expected to outperform both ViT-B and ViT-H, while SAM2 is expected to show a different scaling pattern across Hiera-Small, Hiera-Base+, and Hiera-Large.
        \item[\emph{Baselines:}] SAM ViT-B/L/H and SAM2 Hiera-Small/Base+/Large.
    \end{itemize}

    \item \textbf{RQ2 (Medical pre-training transfer):} Does generic medical pre-training improve zero-shot endoscopic polyp segmentation over vanilla SAM under the same box-prompt protocol?
    \begin{itemize}
        \item[\textbf{$H_2$:}] Generic medical pre-training alone does not guarantee better zero-shot transfer to colonoscopic polyps, because endoscopy differs substantially from the mixed-modality medical data used to train MedSAM.
        \item[\emph{Baselines:}] SAM ViT-B and MedSAM.
    \end{itemize}

    \item \textbf{RQ3 (K-shot adaptation strategy):} Under the same v3 box-prompt protocol, how do full-encoder prompt-perturbation fine-tuning and LoRA-based parameter-efficient fine-tuning compare as the support set increases from 1 to 5 to 10 annotated polyps?
    \begin{itemize}
        \item[\textbf{$H_3$:}] Full-encoder prompt-perturbation fine-tuning is expected to achieve stronger cross-domain Dice than LoRA-only adaptation at the same support-set size, but LoRA is expected to provide a lower-parameter adaptation baseline that directly addresses the Adapter/LoRA requirement in the project brief.
        \item[\emph{Baselines:}] PP-SAM 1/5/10-shot and SAMed LoRA 1/5/10-shot. Both PP-SAM and SAMed LoRA have completed 1/5/10-shot evaluations under the v3 protocol.
    \end{itemize}

    \item \textbf{RQ4 (Gap to full supervision):} How do limited-data adaptation and full-data parameter-efficient adaptation compare with a fully supervised upper-bound reference?
    \begin{itemize}
        \item[\textbf{$H_4$:}] Few-shot support can recover a substantial fraction of the fully supervised performance, but it remains data-limited. SAMed LoRA full 880 is therefore interpreted as a full-data PEFT bridge: it uses the full Kvasir training split but updates only a parameter-efficient subset. Because the current ASPS checkpoint follows the PraNet split rather than the v3 split, ASPS should be interpreted as an external reference rather than a directly ranked v3 baseline.
        \item[\emph{Baselines:}] PP-SAM 5/10-shot for limited-data adaptation, SAMed LoRA full 880 for full-data PEFT, and ASPS as an external fully supervised upper-bound reference.
    \end{itemize}

    \item \textbf{RQ5 (Prompt-free and robustness extensions):} What remains to be resolved for prompt-free, scale-aware, and boundary-robust polyp segmentation beyond the current-stage benchmark?
    \begin{itemize}
        \item[\textbf{$H_5$:}] The current-stage experiments are sufficient to establish a reproducible benchmark and a few-shot adaptation baseline, but they are not sufficient to close the long-term project directions related to full OP-SAM reproduction, systematic prompt-quality perturbation, object-scale robustness, explicit weak-boundary refinement, v3-aligned ASPS retraining, and polyp-specific Medical SAM Adapter training.
        \item[\emph{Baselines and plans:}] The current OP-SAM result is a CPG-only exploratory pilot; full CPG+SPF+EPE reproduction, ASPS retraining on Kvasir \texttt{train.txt}, polyp-specific Medical SAM Adapter training, multi-level prompt perturbation, area-stratified evaluation, and boundary-specific modules are reserved for the next stage.
    \end{itemize}
\end{itemize}

\section{Contributions}
\label{sec:contributions}

The contributions of this partial thesis are four-fold:
\begin{itemize}
    \item A \emph{reproducible v3 evaluation protocol} for promptable polyp segmentation, using Kvasir-SEG \texttt{val.txt} as the held-out in-domain split, CVC-ClinicDB as the cross-domain split, a fixed $\pm30$ pixel box-prompt perturbation, and five metrics: Dice, IoU, precision, recall, and HD95.

    \item A \emph{zero-shot foundation-model benchmark} covering SAM v1 ViT-B/L/H, SAM2 Hiera-Small/Base+/Large, and MedSAM under the same v3 protocol. This benchmark separates backbone-scaling effects from generic medical pre-training effects, rather than attributing all performance differences to adaptation alone.

    \item A \emph{few-shot adaptation comparison} centred on PP-SAM and SAMed. PP-SAM 1/5/10-shot experiments provide the current full-encoder prompt-perturbation fine-tuning baseline, while SAMed LoRA is included as the parameter-efficient Adapter/LoRA extension required by the project brief. At this stage, the SAMed 1/5/10-shot and full 880 runs are available and archived in the formal results directory.

    \item A \emph{transparent separation of aligned results, external references, and future work}. ASPS is reported as an external fully supervised upper-bound reference because the available checkpoint is not v3-aligned, and OP-SAM is reported only as a CPG-only exploratory training-free pilot. Full ASPS-v3 retraining, full OP-SAM reproduction, polyp-specific Medical SAM Adapter training, prompt-quality ablations, object-scale analysis, and weak-boundary refinement are explicitly listed as next-stage objectives rather than claimed as completed work.
\end{itemize}

\section{Organisation of the thesis}
\label{sec:organisation}

The remainder of the thesis is organised as follows. Chapter~2 surveys promptable foundation models, medical SAM adaptations, polyp-specific methods, prompt engineering, datasets, and the long-term gaps that motivate the current-stage benchmark. Chapter~\ref{chap:outline} describes the v3 methodology, including the dataset splits, prompt-generation rule, baseline taxonomy, evaluation metrics, protocol caveats, and the PP-SAM/SAMed adaptation recipes. Chapter~\ref{chap:solution} explains how the methodology is implemented as a reproducible benchmark system, including baseline wrappers, training artifacts, and special-case handling. Chapter~\ref{chap:validation} validates the benchmark through quantitative and qualitative results organised by the research questions. Chapter~\ref{chap:conclusion} concludes the partial thesis by summarising the completed empirical findings and listing the next-stage experiments: ASPS-v3 retraining, full OP-SAM reproduction, polyp-specific Medical SAM Adapter training, prompt-quality perturbation analysis, object-scale analysis, weak-boundary refinement, and extension to additional polyp datasets.
